% concatenated from sn-article.tex
%Version 3.1 December 2024
% See section 11 of the User Manual for version history
%
%%%%%%%%%%%%%%%%%%%%%%%%%%%%%%%%%%%%%%%%%%%%%%%%%%%%%%%%%%%%%%%%%%%%%%
%%                                                                 %%
%% Please do not use \input{...} to include other tex files.       %%
%% Submit your LaTeX manuscript as one .tex document.              %%
%%                                                                 %%
%% All additional figures and files should be attached             %%
%% separately and not embedded in the \TeX\ document itself.       %%
%%                                                                 %%
%%%%%%%%%%%%%%%%%%%%%%%%%%%%%%%%%%%%%%%%%%%%%%%%%%%%%%%%%%%%%%%%%%%%%

%%\documentclass[referee,sn-basic]{sn-jnl}% referee option is meant for double line spacing

%%=======================================================%%
%% to print line numbers in the margin use lineno option %%
%%=======================================================%%

%%\documentclass[lineno,pdflatex,sn-basic]{sn-jnl}% Basic Springer Nature Reference Style/Chemistry Reference Style

%%=========================================================================================%%
%% the documentclass is set to pdflatex as default. You can delete it if not appropriate.  %%
%%=========================================================================================%%

%%\documentclass[sn-basic]{sn-jnl}% Basic Springer Nature Reference Style/Chemistry Reference Style

%%Note: the following reference styles support Namedate and Numbered referencing. By default the style follows the most common style. To switch between the options you can add or remove “Numbered” in the optional parenthesis. 
%%The option is available for: sn-basic.bst, sn-chicago.bst%  
 
%%\documentclass[pdflatex,sn-nature]{sn-jnl}% Style for submissions to Nature Portfolio journals
%%\documentclass[pdflatex,sn-basic]{sn-jnl}% Basic Springer Nature Reference Style/Chemistry Reference Style
\documentclass[pdflatex,sn-mathphys-num]{sn-jnl}% Math and Physical Sciences Numbered Reference Style
%%\documentclass[pdflatex,sn-mathphys-ay]{sn-jnl}% Math and Physical Sciences Author Year Reference Style
%%\documentclass[pdflatex,sn-aps]{sn-jnl}% American Physical Society (APS) Reference Style
%%\documentclass[pdflatex,sn-vancouver-num]{sn-jnl}% Vancouver Numbered Reference Style
%%\documentclass[pdflatex,sn-vancouver-ay]{sn-jnl}% Vancouver Author Year Reference Style
%%\documentclass[pdflatex,sn-apa]{sn-jnl}% APA Reference Style
%%\documentclass[pdflatex,sn-chicago]{sn-jnl}% Chicago-based Humanities Reference Style

%%%% Standard Packages
%%<additional latex packages if required can be included here>
\usepackage{comment}
\usepackage{graphicx}%
\usepackage{multirow}%
\usepackage{amsmath,amssymb,amsfonts}%
\usepackage{amsthm}%
\usepackage{mathrsfs}%
\usepackage[title]{appendix}%
\usepackage{xcolor}%
\usepackage{textcomp}%
\usepackage{manyfoot}%
\usepackage{booktabs}%
\usepackage{algorithm}%
\usepackage{algorithmicx}%
\usepackage{algpseudocode}%
\usepackage{listings}%
%%%%

%%%%%=============================================================================%%%%
%%%%  Remarks: This template is provided to aid authors with the preparation
%%%%  of original research articles intended for submission to journals published 
%%%%  by Springer Nature. The guidance has been prepared in partnership with 
%%%%  production teams to conform to Springer Nature technical requirements. 
%%%%  Editorial and presentation requirements differ among journal portfolios and 
%%%%  research disciplines. You may find sections in this template are irrelevant 
%%%%  to your work and are empowered to omit any such section if allowed by the 
%%%%  journal you intend to submit to. The submission guidelines and policies 
%%%%  of the journal take precedence. A detailed User Manual is available in the 
%%%%  template package for technical guidance.
%%%%%=============================================================================%%%%

%% as per the requirement new theorem styles can be included as shown below
\theoremstyle{thmstyleone}%
\newtheorem{theorem}{Theorem}%  meant for continuous numbers
%%\newtheorem{theorem}{Theorem}[section]% meant for sectionwise numbers
%% optional argument [theorem] produces theorem numbering sequence instead of independent numbers for Proposition
\newtheorem{proposition}[theorem]{Proposition}%
\newtheorem{definition}{Definition}[section]
\newtheorem{example}{Example}[section]
\newtheorem{remark}{Remark}[section]
\newtheorem{corollary}{Corollary}[section]

%%\newtheorem{proposition}{Proposition}% to get separate numbers for theorem and proposition etc.

\theoremstyle{thmstyletwo}%


\theoremstyle{thmstylethree}%


\raggedbottom
%%\unnumbered% uncomment this for unnumbered level heads

\begin{document}

\title[On the Cohomology of Cyclic Associative Algebras]{On the Cohomology of Cyclic Associative Algebras}



\author*[1,2,3]{\fnm{Hassan} \sur{Alhussein}}\email{hassanalhussein2014@gmail.com}



\affil*[1]{ \orgname{Siberian State University of Telecommunication and Informatics Science}, \orgaddress{ \city{Novosibirsk}, \country{Russia}}}

\affil[2]{ \orgname{Novosibirsk State University of Economics and Management}, \orgaddress{ \city{Novosibirsk}, \country{Russia}}}


\affil[3]{ \orgname{Novosibirsk State University}, \orgaddress{ \city{Novosibirsk}, \country{Russia}}}


\abstract{We introduce a cohomology theory for cyclic associative algebras, a subclass of shift associative algebras defined by the identity $(xy)z = x(yz) = y(zx)$. This cohomology, denoted $H^\bullet_{\mathrm{cyc}}(A, M)$, is a subtheory of Hochschild cohomology obtained by restricting to cochains that satisfy a cyclic compatibility condition derived from the defining identity. We prove that $H^2_{\mathrm{cyc}}(A, M)$ classifies cyclic associative extensions of $A$ by a cyclic bimodule $M$. The universal derivation and the module of differential forms $\Omega^\bullet_{\mathbb{F}}(A)$ are constructed, and $(\Omega^\bullet_{\mathbb{F}}(A), d)$ is shown to be the universal cyclic differential graded algebra over $A$. For trivial coefficients, we establish natural inclusions $HC^n(A) \hookrightarrow H^n_{\mathrm{cyc}}(A, \mathbb{F}) \hookrightarrow HH^n(A, \mathbb{F})$, placing our theory intermediate between Connes' cyclic cohomology and Hochschild cohomology.}



\keywords{Shift associative algebras, Cyclic associative algebras, Hochschild cohomology, Cyclic cohomology, Extensions of algebras, Deformations of algebras.}

\pacs[MSC Classification]{17A30, 17D25, 16E40, 18M60, 18G15.}


\maketitle

\section{Introduction}\label{sec1}

The study of non-associative algebras satisfying identities of associative type has a long and rich history. An algebra is said to be of \textit{associative type} if it, or its opposite, satisfies an identity of the form $(x_1x_2)x_3 = x_{\sigma(1)}(x_{\sigma(2)}x_{\sigma(3)})$ or $(x_1x_2)x_3 = (x_{\sigma(1)}x_{\sigma(2)})x_{\sigma(3)}$ for some permutation $\sigma \in \mathbb{S}_3$ \cite{Abdelwahab2025}. This broad class encompasses many important varieties, including associative algebras, commutative associative algebras, Novikov algebras \cite{Dotsenko2023, Sartayev2023}, bicommutative algebras \cite{Drensky2024, Ismailov2024}, perm algebras \cite{Kaygorodov2024, Mashurov2024}, and others. Among these, the variety defined by the identity $(xy)z = y(zx)$ was recently studied by Abdelwahab, Kaygorodov, and Sartayev \cite{Abdelwahab2025} under the name \textit{shift associative algebras}.

A particularly interesting subclass of shift associative algebras consists of those satisfying the stronger condition $(xy)z = x(yz) = y(zx)$. Following \cite{Abdelwahab2025}, we call these \textit{cyclic associative algebras}. This identity imposes a full cyclic symmetry on products of three elements. Examples range from commutative associative algebras to non-commutative nilpotent algebras appearing already in dimension three. The algebraic and geometric classifications of low-dimensional shift associative algebras were established in \cite{Abdelwahab2025}, where it was shown that the first non-associative shift associative algebra appears in dimension five.

The present paper is devoted to developing a cohomology theory for cyclic associative algebras. Cohomology theories for varieties of non-associative algebras play a crucial role in the study of extensions, deformations, and structural properties. For associative algebras, Hochschild cohomology \cite{Hochschild1945} classifies extensions and governs deformation theory. For commutative algebras, Harrison cohomology and André-Quillen cohomology provide analogous tools. For cyclic associative algebras, we introduce a cohomology theory that interpolates between Hochschild cohomology and Connes' cyclic cohomology \cite{Connes1985}.

Our first main result is the construction of a cochain complex $\mathbb{B}^\bullet(A, M)$ for a cyclic associative algebra $A$ with coefficients in a cyclic bimodule $M$. The cochains are multilinear maps $f: A^{\otimes (n+1)} \to M$ satisfying five compatibility conditions (\ref{6})-(\ref{10}) that encode the cyclic symmetry. The differential is the classical Hochschild coboundary, and we prove that it preserves these conditions. The resulting cohomology $H^\bullet_{\mathrm{cyc}}(A, M)$ is a subtheory of Hochschild cohomology.

We then establish the fundamental interpretation of the second cohomology group:

\begin{center}
\emph{$H^2_{\mathrm{cyc}}(A, M)$ is in canonical bijection with the set of equivalence classes of split abelian extensions of $A$ by $M$.}
\end{center}

This classification theorem (Theorem \ref{thm:classification}) justifies the cohomology as the correct tool for studying extensions in the variety of cyclic associative algebras.

To facilitate computations, we construct the enveloping algebra $E(A)$ of a cyclic associative algebra $A$ (Proposition \ref{prop:representation_module}) and show that representations of $A$ correspond to left modules over $E(A)$. The augmentation ideal $I = \ker(E(A) \to A)$ yields the module of differential forms $\Omega^1_{\mathbb{F}}(A) = I/I^2$, and the universal derivation $d: A \to \Omega^1_{\mathbb{F}}(A)$ satisfies the universality property $\operatorname{Der}(A, M) \cong \operatorname{Hom}_A(\Omega^1_{\mathbb{F}}(A), M)$ (Proposition \ref{prop:universal_derivation}). Higher differential forms $\Omega^n_{\mathbb{F}}(A)$ are defined via exterior powers, and we prove that $(\Omega^\bullet_{\mathbb{F}}(A), d)$ is the universal cyclic differential graded algebra over $A$ (Theorem \ref{thm:universal_cdg}).

A key conceptual contribution of this paper is the comparison with existing cohomology theories. For trivial coefficients $M = \mathbb{F}$, we establish natural inclusions:
\[
HC^n(A) \hookrightarrow H^n_{\mathrm{cyc}}(A, \mathbb{F}) \hookrightarrow HH^n(A, \mathbb{F}),
\]
where $HC^n(A)$ denotes Connes' cyclic cohomology and $HH^n(A, \mathbb{F})$ denotes Hochschild cohomology. Thus $H^\bullet_{\mathrm{cyc}}(A, \mathbb{F})$ occupies an intermediate position between these two classical theories (Proposition \ref{prop:comparison}).

Finally, we discuss the notion of smoothness for cyclic associative algebras. An algebra $A$ is called \textit{almost-free} if every abelian extension of $A$ admits a lifting homomorphism. We prove that $A$ is almost-free if and only if $\Omega^1_{\mathbb{F}}(A)$ is projective if and only if $A$ has cohomological dimension $\le 1$ with respect to $H^\bullet_{\mathrm{cyc}}(A, -)$ (Theorem \ref{thm:almost_free_equivalent}). This mirrors classical results in commutative algebra and non-commutative geometry.

The paper is organized as follows. Section 2 recalls the definition and basic properties of cyclic associative algebras and introduces cyclic bimodules. Section 3 constructs the enveloping algebra, the universal derivation, and the module of differential forms, culminating in the universal cyclic differential graded algebra. Section 4 defines the cohomology $H^\bullet_{\mathrm{cyc}}(A, M)$ and proves the extension classification theorem. Section 5 discusses conjectural smoothness and cohomological dimension. Section 6 compares our cohomology with Connes' cyclic cohomology. An appendix contains the detailed step-by-step proof that the Hochschild coboundary preserves the cyclic conditions.







\section{Preliminaries}

\subsection{Cyclic Associative Algebras}

In this subsection, we recall the definition and examples of cyclic associative algebras, which form a subclass of shift associative algebras. Throughout, let $\mathbb{F}$ be an algebraically closed field of characteristic different from $2$.

\begin{definition}\cite{Abdelwahab2025}
An algebra $(A, \cdot)$ is called \textbf{cyclic associative} if it satisfies the following identities for all $x, y, z \in A$:
\begin{align}
(xy)z &= x(yz), \label{eq:assoc}\\
(xy)z &= y(zx). \label{eq:cyclic}
\end{align}
Equivalently,
\[
(xy)z = x(yz) = y(zx).
\]
\end{definition}

\begin{remark}
Recall from \cite{Abdelwahab2025} that a \textbf{shift associative algebra} is defined by the single identity $(xy)z = y(zx)$. Every cyclic associative algebra is shift associative, but the converse is false in general. The first non-associative shift associative algebra appears in dimension $5$ \cite[Proposition 59]{Abdelwahab2025}, whereas non-associative cyclic associative algebras already appear in dimension $3$ (see Example \ref{ex:a1} below).
\end{remark}
\begin{comment}
\begin{proposition}\label{prop:basic}
Let $A$ be a cyclic associative algebra. Then:
\begin{enumerate}
\item[(1)] \textbf{Associativity.} $A$ is associative, i.e., $(xy)z = x(yz)$ for all $x, y, z \in A$.
\item[(2)] \textbf{Commutator properties.} The commutator $[x, y] = xy - yx$ satisfies the Jacobi identity, so $(A, [\cdot, \cdot])$ is a Lie algebra. Moreover, $[[x, y], z] = 0$ for all $x, y, z \in A$; i.e., the Lie algebra is $2$-step nilpotent.
\item[(3)] \textbf{Jordan structure.} The symmetrized product $x \circ y = \frac{1}{2}(xy + yx)$ makes $(A, \circ)$ a commutative associative algebra. In particular, $A$ is power-associative.
\item[(4)] \textbf{4-nice property.} For any four elements $x_1, x_2, x_3, x_4 \in A$, the product $x_1 x_2 x_3 x_4$ is independent of bracketing and order up to cyclic permutation. More precisely, $A$ is $4$-nice in the sense of \cite[Theorem 13]{Abdelwahab2025}.
\item[(5)] \textbf{Nilpotency of the commutator.} The Lie algebra $(A, [\cdot, \cdot])$ is $3$-step nilpotent. If the commutator is trivial, then $A$ is commutative associative.
\item[(6)] \textbf{Idempotents.} Any idempotent $e \neq 0$ in a cyclic associative algebra satisfies $ex = xe = x$ for all $x \in A$ when $A$ is unital; otherwise, the Peirce decomposition simplifies to a direct sum of ideals \cite[Theorem 43]{Abdelwahab2025}.
\end{enumerate}
\end{proposition}

\begin{proof}
Properties (1)--(4) follow directly from the definition and from \cite[Theorem 6, Lemma 9, Theorem 13, Corollary 15]{Abdelwahab2025}. Property (5) is \cite[Corollary 15]{Abdelwahab2025}. Property (6) is \cite[Corollary 24 and Theorem 43]{Abdelwahab2025}. ∎
\end{proof}

We now present examples of cyclic associative algebras, ranging from commutative to genuinely non-commutative cases.
\end{comment}
\begin{example}\label{ex:comm}
 Every commutative associative algebra is cyclic associative. 
\end{example}

\begin{example}\label{ex:a1}
 Let $A$ have basis $\{e_1, e_2,e_3\}$ with multiplication:
\[
e_1e_2 = e_3,\quad e_2e_1 = -e_3,
\]
all other products zero. This algebra is cyclic associative. It is non-commutative. This is the smallest non-commutative cyclic associative algebra.
\end{example}
\begin{example}\label{prop:construction_f}
Let $A$ be a commutative associative algebra and let $M$ be a right $A$-module. Let $f: M \to A$ be an $A$-module morphism satisfying
\[
a \cdot f(b) = b \cdot f(a) \qquad \forall a, b \in M.
\]
Define a product $\ast$ on $M$ by
\[
a \ast b = a \cdot f(b) \in M.
\]
Then $(M, \ast)$ is a cyclic associative algebra.
\end{example}
\begin{comment}
    

\begin{example}\label{ex:a2alpha}
\textbf{The $3$-dimensional algebra $\mathfrak{a}_2^\alpha$ from \cite[Theorem 54]{Abdelwahab2025}.} For any $\alpha \in \mathbb{F}$, define $A$ with basis $\{e_1, e_2, e_3\}$ and multiplication:
\[
e_1e_1 = e_3,\quad e_1e_2 = e_3,\quad e_2e_1 = -e_3,\quad e_2e_2 = \alpha e_3,
\]
all other products zero. This family is cyclic associative. When $\alpha = 0$, we recover a $2$-step nilpotent algebra. These algebras are pairwise non-isomorphic for distinct $\alpha$.
\end{example}

\begin{example}\label{ex:a13}
\textbf{The $4$-dimensional algebra $\mathfrak{a}_{13}$ from \cite[Theorem 57]{Abdelwahab2025}.} Let $A$ have basis $\{e_1, e_2, e_3, e_4\}$ with multiplication:
\[
e_1e_1 = e_4,\quad e_1e_2 = e_3,\quad e_2e_1 = -e_3,\quad e_2e_2 = e_3,\quad e_1e_4 = e_3,\quad e_4e_1 = e_3,
\]
all other products zero. This algebra is cyclic associative but not commutative, and it is not $2$-step nilpotent (since $e_4e_4 = 0$ but products involving $e_4$ create $e_3$).
\end{example}

\begin{example}
\textbf{The unital hull of a nilpotent commutative algebra.} Let $B$ be a nilpotent commutative associative algebra. Then the unital hull $\hat{B} = \mathbb{F}1 \oplus B$ with multiplication
\[
(\alpha 1 + b)(\beta 1 + c) = \alpha\beta 1 + (\alpha c + \beta b + bc)
\]
is cyclic associative and unital. By \cite[Lemma 10]{Abdelwahab2025}, any unital cyclic associative algebra is commutative associative, so this example is commutative.
\end{example}

\begin{remark}
The algebraic classification of all complex cyclic associative algebras of dimensions $\le 4$ is given in \cite[Theorem 57]{Abdelwahab2025}. In dimension $3$, the non-commutative algebras are exactly $\mathfrak{a}_1$ and $\mathfrak{a}_2^\alpha$ above. In dimension $4$, there are $14$ families of non-commutative cyclic associative algebras, including $\mathfrak{a}_{13}$ and $\mathfrak{a}_{14}$ given in \cite{Abdelwahab2025}. We will use this classification in Section 6 to compute cohomology groups explicitly.
\end{remark}

    

\begin{notation}
Throughout the remainder of the paper, we denote:
\begin{itemize}
\item $A$: a cyclic associative algebra.
\item $M$: a cyclic $A$-bimodule (defined in Section 2.2).
\item $A^{\mathrm{ab}} = A / [A, A]$: the abelianization of $A$, which is a commutative associative algebra.
\item $[x, y] = xy - yx$: the commutator.
\item $x \circ y = \frac{1}{2}(xy + yx)$: the Jordan (symmetrized) product.
\end{itemize}
\end{notation}




\bmhead{Acknowledgements}

Acknowledgements are not compulsory. Where included they should be brief. Grant or contribution numbers may be acknowledged.

Please refer to Journal-level guidance for any specific requirements.
\end{comment}





%%===================================================%%
%% For presentation purpose, we have included        %%
%% \bigskip command. Please ignore this.             %%
%%===================================================%%


\subsection{Cyclic Bimodules}

Let $A$ be a cyclic associative algebra. We now introduce the notion of a bimodule over $A$ that is compatible with the cyclic identity. The guiding principle is that the defining identities of $A$ should extend to any expression involving elements of $A$ and a module element $m \in M$, with $m$ placed in any position.

\begin{definition}\label{def:cyclic_bimodule}
Let $A$ be a cyclic associative algebra. A \textbf{cyclic $A$-bimodule} is a vector space $M$ equipped with two bilinear maps
\[
A \otimes M \to M,\quad (a,m) \mapsto a \cdot m,
\qquad
M \otimes A \to M,\quad (m,a) \mapsto m \cdot a,
\]
such that for all $x, y, z \in A$ and $m \in M$, the following identities hold:
\begin{align}
\text{(C1)}&\quad (m x) y = m (x y) = x (y m), \label{eq:cyc1}\\
\text{(C2)}&\quad (x m) y = x (m y) = m (y x), \label{eq:cyc2}\\
\text{(C3)}&\quad (x y) m = x (y m) = y (m x). \label{eq:cyc3}
\end{align}
\end{definition}
\subsection{Enveloping Algebra of a Cyclic Associative Algebra}

Let $A_l$ and $A_r$ be two copies of $A$, denoted by
\[
A_l = A\otimes \mathbb{F}, \qquad A_r = \mathbb{F}\otimes A.
\]
Consider the $\mathbb{F}$-module
\[
E(A) = (A \otimes A) \oplus A_l \oplus A_r.
\]

Define an associative product on $E(A)$ by the following rules, extended bilinearly:
\[
\begin{aligned}
(a \otimes b)(c \otimes d) &= a \otimes (b c d), \\
(a \otimes b) c_l &= a \otimes (b c), \\
c_l (a \otimes b) &= (c a) \otimes b, \\
(a \otimes b) c_r &= a \otimes (b c), \\
c_r (a \otimes b) &= a \otimes (c b), \\
a_l b_l &= (ab)_l, \\
a_r b_r &= (ba)_r, \\
a_l b_r &= (ab)_r, \\
a_r b_l &= (ba)_l,
\end{aligned}
\]
and all other products are zero.

\begin{proposition}\label{prop:representation_module}
Every representation $M$ of a cyclic associative algebra $A$ is equivalent to a left module over the enveloping algebra $E(A)$ via
\[
(a \otimes b) \cdot m = a \cdot (m \cdot b), \quad a_l \cdot m = a \cdot m, \quad a_r \cdot m = m \cdot a.
\]
Conversely, any left $E(A)$-module $M$ gives a representation of $A$ by $a \cdot m = a_l \cdot m$ and $m \cdot a = a_r \cdot m$.
\end{proposition}

\begin{proof}
The verification follows directly from the cyclic associative identities of $A$ and the cyclic bimodule conditions (C1)--(C3). 
\end{proof}

\begin{proposition}\label{prop:augmentation}
The map $\mu: E(A) \to A$ defined by
\[
\mu(a \otimes b) = ab, \qquad \mu(a_l) = a, \qquad \mu(a_r) = a,
\]
is a surjective algebra homomorphism. Let $I = \ker \mu$. The quotient
\[
\Omega^1_{\mathbb{F}}(A) = I / I^2
\]
is called the \textbf{module of differential forms} of $A$.
\end{proposition}

\begin{proof}
Surjectivity is clear. For multiplicativity, one checks each case using the associativity of $A$ and the cyclic identities. For instance,
\[
\mu((a \otimes b)(c \otimes d)) = \mu(a \otimes (b c d)) = a(b c d),
\]
\[
\mu(a \otimes b)\mu(c \otimes d) = (ab)(cd) = a(b(cd)) = a(b c d),
\]
where the last equality uses associativity. The other cases are similar. 
\end{proof}
\subsection{Universal Derivation of a Cyclic Associative Algebra}

Let $A$ be a cyclic associative algebra and let $M$ be a representation of $A$ (a cyclic $A$-bimodule).

\begin{definition}\label{def:derivation}
A \textbf{derivation} from $A$ to $M$ is a $\mathbb{F}$-linear map $D: A \to M$ satisfying the Leibniz rule:
\[
D(ab) = D(a) \cdot b + a \cdot D(b), \qquad \forall a, b \in A.
\]
\end{definition}
We denote by $Der(A,M)$ the set of all derivations from $A$ to $M$.
\begin{definition}\label{def:inner_derivation}
For a fixed element $m \in M$, the map
\[
\operatorname{ad}_m: A \to M, \qquad \operatorname{ad}_m(a) = a \cdot m - m \cdot a,
\]
is called an \textbf{inner derivation}. The set of all inner derivations is denoted by $\operatorname{Inn}(A, M)$.
\end{definition}

\begin{definition}\label{def:universal_derivation}
A derivation $d: A \to \Omega^1_{\mathbb{F}}(A)$ is called \textbf{universal} if for any derivation $D: A \to M$ into a representation $M$, there exists a unique $A$-linear map $\phi: \Omega^1_{\mathbb{F}}(A) \to M$ such that $D = \phi \circ d$.
\end{definition}

\begin{proposition}\label{prop:universal_derivation}
Let $A$ be a cyclic associative algebra and let $E(A)$ be its enveloping algebra with augmentation ideal $I = \ker(\mu: E(A) \to A)$. Then
\[
\Omega^1_{\mathbb{F}}(A) = I / I^2
\]
is a representation of $A$, and the map
\[
d: A \to \Omega^1_{\mathbb{F}}(A), \qquad d(a) = a_l - a_r \pmod{I^2},
\]
is a universal derivation.
\end{proposition}

\begin{proof}
We first show that $d$ is a derivation. For any $a, b \in A$, we have in $E(A)$:
\[
d(ab) = (ab)_l - (ab)_r = a_l b_l - a_r b_r.
\]
Adding and subtracting $a_l b_r$ and $a_r b_l$:
\[
d(ab) = (a_l b_l - a_l b_r) + (a_l b_r - a_r b_r) + (a_r b_l - a_r b_r) + (a_l b_r - a_r b_l).
\]
In $I/I^2$, we have $b_l \equiv b_r$ and $a_l \equiv a_r$, so:
\[
d(ab) \equiv (a_l - a_r) b_l + a_r (b_l - b_r) \equiv d(a) \cdot b + a \cdot d(b) \pmod{I^2}.
\]
Hence $d(ab) = d(a) b + a d(b)$ in $\Omega^1_{\mathbb{F}}(A)$.

Now let $D: A \to M$ be any derivation. Define $\phi: \Omega^1_{\mathbb{F}}(A) \to M$ on generators by
\[
\phi(d(a)) = D(a), \qquad \forall a \in A,
\]
and extend by $A$-linearity. We must check that $\phi$ is well-defined, i.e., that it vanishes on $I^2$. Since $I$ is generated by elements of the form $a_l - a_r$ and $a \otimes b - (ab)_r$, it suffices to verify:
\[
\phi((a_l - a_r) b) = D(a) b, \qquad \phi(a (b_l - b_r)) = a D(b).
\]
But these follow from the derivation property of $D$ and the definition of $\phi$. Hence $\phi$ is well-defined. Uniqueness follows from the fact that the elements $d(a)$ generate $\Omega^1_{\mathbb{F}}(A)$ as an $A$-module. 
\end{proof}

\begin{corollary}\label{cor:der_hom}
For any representation $M$ of a cyclic associative algebra $A$, there is a natural isomorphism
\[
\operatorname{Der}(A, M) \cong \operatorname{Hom}_A(\Omega^1_{\mathbb{F}}(A), M).
\]
\end{corollary}

\subsection{Differential Forms}

\begin{definition}\label{def:differential_forms}
Let $A$ be a cyclic associative algebra. For $n \ge 1$, define the \textbf{module of $n$-differential forms} by
\[
\Omega^n_{\mathbb{F}}(A) = \bigwedge^n_A \Omega^1_{\mathbb{F}}(A),
\]
where $\bigwedge_A$ denotes the exterior product over $A$. More explicitly,
\[
\Omega^n_{\mathbb{F}}(A) = \Omega^1_{\mathbb{F}}(A) \otimes_A \bigwedge_A^{n-1} \Omega^1_{\mathbb{F}}(A).
\]
We set $\Omega^0_{\mathbb{F}}(A) = A$.
\end{definition}

The space $\Omega^n_{\mathbb{F}}(A)$ is generated by symbols
\[
\omega = a_0 \, da_1 \otimes da_2 \wedge \cdots \wedge da_n + db_1 \otimes db_2 \wedge \cdots \wedge db_n, \qquad a_i, b_j \in A.
\]

\begin{definition}\label{def:differential_operator}
Define a linear map $d: \Omega^n_{\mathbb{F}}(A) \to \Omega^{n+1}_{\mathbb{F}}(A)$ by
\[
d(a_0 \, da_1 \otimes da_2 \wedge \cdots \wedge da_n) = da_0 \otimes da_1 \wedge da_2 \wedge \cdots \wedge da_n,
\]
and extend by linearity. For $n = 0$, we set $d(a) = da$ as defined in Proposition \ref{prop:universal_derivation}.
\end{definition}

\begin{proposition}\label{prop:differential_properties}
The map $d: \Omega^\bullet_{\mathbb{F}}(A) \to \Omega^{\bullet+1}_{\mathbb{F}}(A)$ satisfies:
\begin{enumerate}
\item[(1)] $d \circ d = 0$ (i.e., $(\Omega^\bullet_{\mathbb{F}}(A), d)$ is a cochain complex);
\item[(2)] $d(\omega \wedge \eta) = d\omega \wedge \eta + (-1)^{\deg \omega} \omega \wedge d\eta$ (graded Leibniz rule);
\item[(3)] $d(ab) = d(a) b + a d(b)$ for all $a, b \in A$ (extends the derivation property).
\end{enumerate}
\end{proposition}

\begin{proof}
(1) For any generator $a_0 da_1 \wedge \cdots \wedge da_n$, we have
\[
d^2(a_0 da_1 \wedge \cdots \wedge da_n) = d(da_0 \wedge da_1 \wedge \cdots \wedge da_n) = d^2(a_0) \wedge da_1 \wedge \cdots \wedge da_n = 0,
\]
since $d^2(a_0) = 0$ in $\Omega^2_{\mathbb{F}}(A)$.

(2) The graded Leibniz rule follows from the definition of the wedge product and the fact that $d$ is a derivation on $A$.

(3) This is precisely the universal derivation property. 
\end{proof}

\begin{definition}\label{def:cdg_algebra}
A \textbf{cyclic differential graded algebra} (abbreviated \textbf{CDG algebra}) is a graded $\mathbb{F}$-module $P_\bullet = \bigoplus_{n \ge 0} P_n$ equipped with a product $P_r \otimes P_s \to P_{r+s}$ and a differential $d: P_r \to P_{r+1}$ satisfying:
\[
\begin{cases}
d(ab) = (da)b + (-1)^{rs} a(db), \\
d \circ d = 0, \\
a(bc) = b(ca) \quad \text{(cyclic identity in graded sense)}.
\end{cases}
\]
\end{definition}

\begin{theorem}\label{thm:universal_cdg}
The algebra $(\Omega^\bullet_{\mathbb{F}}(A), d)$ is the universal CDG algebra over $A$. That is, for any CDG algebra $P_\bullet$ and any morphism of cyclic associative algebras $\phi: A \to P_0$, there exists a unique morphism of CDG algebras $\widetilde{\phi}: \Omega^\bullet_{\mathbb{F}}(A) \to P_\bullet$ extending $\phi$ and commuting with $d$.
\end{theorem}

\begin{proof}
Define $\widetilde{\phi}$ on generators by
\[
\widetilde{\phi}(a_0 da_1 \wedge \cdots \wedge da_n) = \phi(a_0) d\phi(a_1) \wedge \cdots \wedge d\phi(a_n),
\]
and extend by linearity. Using the universal property of $\Omega^1_{\mathbb{F}}(A)$ and the fact that $P_\bullet$ is a CDG algebra, one verifies that $\widetilde{\phi}$ is well-defined and commutes with $d$. Uniqueness follows from the fact that the elements $a_0 da_1 \wedge \cdots \wedge da_n$ generate $\Omega^\bullet_{\mathbb{F}}(A)$. 
\end{proof}

\begin{corollary}\label{cor:derivative_representation}
For any representation $M$ of $A$, there is a natural isomorphism
\[
\operatorname{Hom}_A(\Omega^n_{\mathbb{F}}(A), M) \cong \operatorname{Alt}^n_A(A, M),
\]
where $\operatorname{Alt}^n_A(A, M)$ denotes the space of skew-symmetric $n$-derivations from $A$ to $M$.
\end{corollary}
\section{Cohomology of Cyclic Associative Algebras}

\subsection{Abelian extensions of a cyclic associative algebra}

Here we try to define the cohomology of cyclic associative algebras by studying abelian extensions of a given cyclic associative algebra $A$ by an $A$-representation $M$, fixed for this paragraph.

An \textbf{abelian extension} of $A$ by $M$ is a short exact sequence of cyclic associative algebras
\[
0 \longrightarrow M \xrightarrow{j} E \xrightarrow{p} A \longrightarrow 0
\]
where $M$ is seen as an abelian cyclic associative algebra (i.e., $mm' = 0$, $\forall m, m' \in M$).

Two such extensions $(E)$ and $(E')$ are said to be \textbf{equivalent} if there exists a cyclic associative algebra morphism $\phi: E \to E'$ such that
\[
p' \circ \phi = p \quad \text{and} \quad \phi \circ j = j',
\]
that is, making commutative the diagram:
\[
\begin{array}{ccccccccc}
0 & \longrightarrow & M & \xrightarrow{j} & E & \xrightarrow{p} & A & \longrightarrow & 0 \\
& & \| & & \downarrow{\phi} & & \| & & \\
0 & \longrightarrow & M & \xrightarrow{j'} & E' & \xrightarrow{p'} & A & \longrightarrow & 0.
\end{array}
\]

By the Five Lemma, such a morphism $\phi$ is necessarily bijective.

For instance, the direct sum $A \oplus M$ is an abelian extension of $A$ by $M$ (with the obvious inclusion and projection). Moreover, this latter is trivially \textbf{split}. In fact, one says that an abelian extension $(E)$ is \textbf{split} if there exists a $\mathbb{F}$-linear map $s: A \to E$ such that $p \circ s = \operatorname{id}_A$. Moreover, if the section $s$ is a cyclic associative algebra morphism, then the extension $(E)$ is said to be \textbf{strongly split} or \textbf{inessential}.

Any abelian extension $(E)$ with a section $s$ yields another $A$-representation structure on $M$ by
\[
a \cdot m = s(a) j(m) \quad \text{and} \quad m \cdot a = j(m) s(a), \qquad \forall a \in A, \, m \in M,
\]
the last products being taken in $E$ (this has a sense since $M = \ker p$). This new structure is naturally independent of the choice of the section $s$. Indeed, if $s'$ is another section of $p$, then we have
\[
p(s'(a) - s(a)) = p s'(a) - p s(a) = a - a = 0,
\]
that is, $s'(a) - s(a) \in \ker p = \operatorname{im} j$. And since $M$ is abelian, we have
\[
s'(a) j(m) = s(a) j(m) \quad \text{and} \quad j(m) s'(a) = j(m) s(a), \qquad \forall m \in M, \, \forall a \in A.
\]

From now on, we are interested only in the set of equivalence classes of split abelian extensions such that the $A$-representation structure of $M$ is the prescribed one.

\subsection{Construction of split abelian extensions}

Let us consider a $\mathbb{F}$-bilinear map $f: A \times A \to M$ and let $A \oplus_f M$ be the $\mathbb{F}$-module $A \oplus M$ equipped with the product given by
\begin{equation}\label{1}
(a, m) \cdot (b, m') = (ab, \; a m' + m b + f(a, b)). 
\end{equation}

Then it is straightforward to check that the product \ref{1} is associative if and only if
\begin{equation}\label{2}
a f(b, c) + f(a b, c) = f(a, b c) + f(a, b) c, \qquad \forall a, b, c \in A,
\end{equation}
and is cyclic associative (i.e., satisfies $(xy)z = y(zx)$) if and only if
\begin{equation}\label{3}
a f(b, c) + f(a, b c) = b f(c, a) + f(b, c a), \qquad \forall a, b, c \in A. 
\end{equation}

In fact, relation \ref{2} can be rewritten as
\[
a f(b, c) - f(ab,c) + f(a,bc) - f(a, b) c = 0,
\]
which is nothing but the $2$-cocyclicity condition for the Hochschild coboundary of an associative algebra. On the other hand, relation \ref{3} expresses the cyclic condition:
\[
a f(b, c) + f(a, bc) = b f(c, a) + f(b, ca).
\]
It can be seen as conjection of the following generlatilaztion 



For all $a, a_1, \ldots, a_{n+1} \in A$ and any $i = 1, \ldots, n+1$,
\begin{equation}\label{6}
a \cdot f(a_1, \ldots, a_i, \ldots, a_{n+1}) = a_i \cdot f(a_{i+1}, \ldots, a_{n+1},a, a_1, \ldots, a_{i-1}). 
\end{equation}


For $j = 1, \ldots, n$,
\begin{equation}\label{7}
f(a_0, a_1, \ldots, a_{j-1}, b a_j, a_{j+1}, \ldots, a_n) = f(a_0, a_1, \ldots, b, a_j a_{j-1}, a_{j+1}, \ldots, a_n). 
\end{equation}

\begin{equation}\label{8}
a \cdot f(b a_0, a_1, \ldots, a_n) = b \cdot f(a_0 a, a_1, \ldots, a_n). 
\end{equation}

For $j = 0, \ldots, n$,
\begin{equation}\label{9}
a \cdot a'_j \cdot f(a_0, a_1, \ldots, a_j, \ldots, a_n) = a'_j \cdot a_j \cdot f(a_0, a_1, \ldots, a, \ldots, a_n). 
\end{equation}


For $i = 1, \ldots, n+1$,
\begin{equation}\label{10}
a \cdot f(a_1, \ldots, a_i b, \ldots, a_{n+1}) = a_i\cdot b  \cdot f(a_{i+1}, \ldots, a_{n+1},a, a_1, \ldots, a_{i-1}). 
\end{equation}
\begin{proposition}\label{prop:cohomology_preserved}
If $f: A^{\otimes (n+1)} \to M$ satisfies \ref{6}-\ref{10}, then so does the Hochschild coboundary $\delta(f)$.
\end{proposition}
\begin{proof}[Sketch]
The verification is a direct computation using the cyclic associative identities of $A$ and the cyclic bimodule conditions. The full step-by-step proof is provided in Appendix \ref{appendix:proof}.
\end{proof}
\subsection{Cohomology of a cyclic associative algebra}

According to Proposition \ref{prop:cohomology_preserved}, we define the \textbf{cohomology of a cyclic associative algebra} $A$ with values in an $A$-representation $M$ to be the cohomology of the complex $(\mathbb{B}^\bullet(A, M), \delta)$ where

\[
\begin{aligned}
\mathbb{B}^0(A, M) &:= M, \\
\mathbb{B}^1(A, M) &:= \{ f: A \to M \mid a f(bc) = b f(ca), \ \forall a, b, c \in A \}, \\
\mathbb{B}^n(A, M) &:= \{ f \in \operatorname{Hom}(A^{\otimes (n+1)}, M) \mid f \text{ satisfies } \ref{6}-\ref{10} \}, \quad n \ge 2,
\end{aligned}
\]

and the Hochschild coboundary $\delta$ acts as usual by

\[
\begin{aligned}
(\delta f)(a_0, \ldots, a_{n+1}) &:= a_0 f(a_1, \ldots, a_{n+1}) \\
&\quad + \sum_{i=0}^{n} (-1)^{i+1} f(a_0, \ldots, a_i a_{i+1}, \ldots, a_{n+1}) \\
&\quad + (-1)^{n+1} f(a_0, \ldots, a_n) a_{n+1}.
\end{aligned}
\]

We denote this cohomology by $H^\bullet_{\mathrm{cyc}}(A, M)$ or simply by $H^\bullet_{\mathrm{cyc}}(R) := H^\bullet_{\mathrm{cyc}}(A, A)$.

For $n = 0$, $H^0_{\mathrm{cyc}}(A, M)$ is the submodule of \textbf{invariants} of $M$, that is,
\[
H^0_{\mathrm{cyc}}(A, M) = M^A := \{ m \in M \mid a m = m a, \ \forall a \in A \}.
\]

For $n = 1$, a $1$-cocycle is a $\mathbb{F}$-linear map $D: A \to M$ such that
\[
D(ab) = D(a) b + a D(b), \quad \forall a, b \in A,
\]
that is, a \textbf{derivation} from $A$ to $M$. Observe that the additional relation $a D(bc) = b D(ca)$ is fulfilled by any derivation thanks to the cyclic of $M$ (i.e., $m \cdot (ab) = b \cdot (ma)$). The set of all derivations from $A$ to $M$ is denoted by $\operatorname{Der}(A, M)$ and we have
\[
H^1_{\mathrm{cyc}}(A, M) = \operatorname{Der}(A, M) / \operatorname{Inn}(A, M),
\]
where $\operatorname{Inn}(R, M)$ is the subset of \textbf{inner derivations}, i.e., derivations of the form $\operatorname{ad}_m(r) = [m, r] = m r - r m$ for a fixed element $m \in M$.

For $n = 2$, we have the following classical classification theorem.

\begin{theorem}\label{thm:classification}
Let $A$ be a cyclic associative algebra and let $M$ be an $A$-representation. Then, there is a canonical bijection
\[
H^2_{\mathrm{cyc}}(A, M) \cong \mathbf{Ext}(A, M),
\]
that is, the set of equivalence classes of split abelian extensions of $A$ by $M$.
\end{theorem}

\begin{proof}
By construction of the cohomology $H^\bullet_{\mathrm{cyc}}(A, M)$, any element $f \in \mathbb{B}^2(A, M)$ is a $2$-cocycle if and only if the algebra $A \oplus_f M$ defined by (5.1) is a cyclic associative algebra. To be more precise, any $2$-cocycle $f \in \mathbb{B}^2(A, M)$ determines a split abelian extension of $A$ by $M$, and any split abelian extension defines a $2$-cocycle $f \in \mathbb{B}^2(A, M)$ by
\[
f(a, b) = s(ab) - s(a) s(b), \quad \forall a, b \in R,
\]
where $s: A \to E$ is a section splitting the extension. The cocycle $f$ measures the obstruction for $s$ to be a cyclic associative algebra morphism, that is, the obstruction to the extension being inessential.

Therefore, it is left to us to show that its equivalence class, characterized by the morphism $\phi$, only involves coboundaries of $\mathbb{B}^1(A, M)$. To this end, let $\phi: A \oplus_f M \to E$ be an equivalence of abelian extensions, i.e., a commutative diagram
\[
\begin{array}{ccccccccc}
0 & \longrightarrow & M & \xrightarrow{j} & A \oplus_f M & \xrightarrow{\pi} & A & \longrightarrow & 0 \\
& & \| & & \downarrow{\phi} & & \| & & \\
0 & \longrightarrow & M & \xrightarrow{j'} & E & \xrightarrow{p} & A & \longrightarrow & 0.
\end{array}
\]
Denoting by $\sigma: R \to R \oplus_f M$, $r \mapsto (r, 0)$, the trivial section of $\pi$, we have
\[
p \circ (\phi \circ \sigma) = (p \circ \phi) \circ \sigma = \pi \circ \sigma = \operatorname{id}_R.
\]
Therefore the map $s := \phi \circ \sigma$ is also a section of $p$. It corresponds to a $2$-cocycle $f' \in \mathbb{B}^2(R, M)$ related to the extension $E$ by the same formula. Since the map $f'$ is given by $f'(a, b) = s(ab) - s(a) s(b)$, one easily checks that the difference $f - f'$ is nothing but the coboundary $\delta(g)$ where $g: R \to M$ is the map $a \mapsto g(a) := \sigma(a) - s(a)$. In fact, a priori, the map $g$ takes its values in the direct sum $R \oplus M$. But since the initial $R$-representation structure of $M$ coincides with the structure induced by the sections, we have $\pi \circ g = p \circ g = 0$. So $\operatorname{im}(g) \subset M$; from whence the theorem. 
\end{proof}
\section{Conjectural smoothness}

From now on, we refer to the paper by Cuntz and Quillen \cite{CuntzQuillen}. Let $R$ be a cyclic associative algebra with an abelian ideal $M$ and let $p: R \to A$ be a surjective morphism such that the sequence
\[
0 \longrightarrow M \xrightarrow{j} R \xrightarrow{p} A \longrightarrow 0
\]
is exact. We are looking for a cyclic associative algebra morphism $l: A \to R$ such that $p \circ l = \operatorname{id}_A$. Then we have an isomorphism $R \cong A \oplus M$ relative to which $l$ becomes an inclusion of $A$. Therefore $H^2_{\mathrm{cyc}}(A, M) \cong \mathbf{Ext}(A, M) = 0$ for all $R$-representations $M$.

\begin{definition}\label{def:almost_free}
A cyclic associative algebra $A$ is called \textbf{almost-free} when for any abelian extension $R$ of $A$, there exists a lifting homomorphism $A \to R$.
\end{definition}

We expect an interpretation of the cohomology theory $H^\bullet_{\mathrm{cyc}}(A, M)$ as $\mathbf{Ext}_{E(A)}^\bullet(A, M)$, and the exact sequence
\[
0 \longrightarrow \Omega^1_{\mathbb{F}}(A) \longrightarrow E(A) \longrightarrow A \longrightarrow 0
\]
could yield the isomorphism
\[
H^{n+1}_{\mathrm{cyc}}(A, M) \cong \mathbf{Ext}_{E(A)}^{n+1}(A, M) \cong \mathbf{Ext}_{E(A)}^n(\Omega^1_{\mathbb{F}}(A), M).
\]
These suppose the construction of derived functors in a non-unitary context.

\begin{theorem}\label{thm:almost_free_equivalent}
Let $A$ be a cyclic associative algebra. The following conditions are equivalent:
\begin{enumerate}
\item[(1)] $A$ is almost-free;
\item[(2)] the $A$-bimodule $\Omega^1_{\mathbb{F}}(A)$ is projective;
\item[(3)] $A$ has cohomological dimension $\le 1$ with respect to the cyclic associative cohomology $H^\bullet_{\mathrm{cyc}}(A, -)$.
\end{enumerate}
\end{theorem}

\begin{proof}
The proof follows the same lines as in the classical case of associative algebras, adapted to the cyclic setting. The key point is that the universal derivation $d: A \to \Omega^1_{\mathbb{F}}(A)$ and the exact sequence
\[
0 \longrightarrow \Omega^1_{\mathbb{F}}(A) \longrightarrow E(A) \longrightarrow A \longrightarrow 0
\]
play the same role as in the Hochschild cohomology theory. The projectiveivity of $\Omega^1_{\mathbb{F}}(A)$ implies that any derivation $D: A \to M$ factors through a projective module, which allows lifting of homomorphisms. Conversely, if $A$ is almost-free, then the extension corresponding to $\Omega^1_{\mathbb{F}}(A)$ splits, so $\Omega^1_{\mathbb{F}}(A)$ is projective. The equivalence with cohomological dimension $\le 1$ follows from the long exact sequence of cohomology associated to the short exact sequence above. 
\end{proof}

We can see that any classical smooth algebra $A$ (commutative and unital) is almost-free. Also, any free cyclic associative algebra is almost-free.

\begin{theorem}\label{thm:cohomological_dimension_zero}
Let $A$ be a cyclic associative algebra. The following properties are equivalent:
\begin{enumerate}
\item[(1)] $A$ has cohomological dimension zero with respect to cyclic associative cohomology, i.e., $H^1_{\mathrm{cyc}}(A, M) = 0$ for all representations $M$;
\item[(2)] the $A$-module $A$ is projective;
\item[(3)] any derivation $D: A \to M$ is inner.
\end{enumerate}
\end{theorem}

\begin{proof}
The equivalence between (1) and (3) follows from the interpretation of $H^1_{\mathrm{cyc}}(A, M)$ as derivations modulo inner derivations. If $H^1_{\mathrm{cyc}}(A, M) = 0$ for all $M$, then every derivation is inner. In particular, the universal derivation $d: A \to \Omega^1_{\mathbb{F}}(A)$ is inner, so $\Omega^1_{\mathbb{F}}(A)$ is a direct summand of $A$, hence projective. The converse is clear. 
\end{proof}
\section{Relation to Connes' Cyclic Cohomology}

Let $A$ be a cyclic associative algebra over a field $\mathbb{F}$ of characteristic $0$. We compare the cohomology theory $H^\bullet_{\mathrm{cyc}}(A, \mathbb{F})$ (with trivial coefficients) with Connes' cyclic cohomology $HC^\bullet(A)$.

\subsection{Connes' Cyclic Cohomology (Brief Recall)}

For an associative algebra $A$, Connes' cyclic cohomology $HC^n(A)$ is the cohomology of the total complex of the $(b, B)$-bicomplex, where $b$ is the Hochschild differential and $B$ is Connes' boundary operator. Equivalently, $HC^n(A) = H^n(C^\bullet_\lambda(A))$ with
\[
C^n_\lambda(A) = \operatorname{Hom}(A^{\otimes (n+1)}, \mathbb{F}) / (1 - \lambda),
\]
and $\lambda$ is the cyclic operator:
\[
(\lambda f)(a_0, \dots, a_n) = (-1)^n f(a_n, a_0, \dots, a_{n-1}).
\]

\subsection{Comparison}

Both cohomology theories incorporate cyclic symmetry, but in different ways:

\begin{itemize}
\item \textbf{Connes' cyclic cohomology} uses the operator $\lambda$ on the cochain complex and the differential $b + B$. It is defined only for trivial coefficients $\mathbb{F}$.
\item \textbf{Cyclic associative cohomology} $H^\bullet_{\mathrm{cyc}}(A, M)$ uses the Hochschild differential $b$ but restricts to cochains that satisfy the cyclic condition (5.4). It is defined for any cyclic bimodule $M$.
\end{itemize}

For trivial coefficients $M = \mathbb{F}$, there is a natural map:
\[
\iota: HC^n(A) \longrightarrow H^n_{\mathrm{cyc}}(A, \mathbb{F}).
\]

Indeed, a Connes cyclic cochain $f \in C^n_\lambda(A)$ satisfies $(1-\lambda)f = 0$, which implies the cyclic conditionS \ref{6}-\ref{10} for the corresponding Hochschild cochain. Moreover, $Bf$ vanishes in $H^n_{\mathrm{cyc}}$ because it is a coboundary.



\begin{proposition}\label{prop:comparison}
Let $A$ be a cyclic associative algebra. There is a natural injective map
\[
HC^n(A) \hookrightarrow H^n_{\mathrm{cyc}}(A, \mathbb{F}) \hookrightarrow HH^n(A, \mathbb{F}),
\]
where $HH^n$ denotes Hochschild cohomology. In general, these inclusions are strict.
\end{proposition}

\begin{proof}
The first map sends a cyclic cochain $f$ to its class in $H^n_{\mathrm{cyc}}$. Injectivity follows from the fact that if $f = (b + B)g$ in Connes' theory, then $f = b g$ in $H^n_{\mathrm{cyc}}$ because $B g$ is a coboundary. The second map is the inclusion of cyclic cochains into all cochains. Examples from low-dimensional classifications show that both maps are not surjective in general. 
\end{proof}

\begin{corollary}
The cyclic associative cohomology $H^\bullet_{\mathrm{cyc}}(A, \mathbb{F})$ sits between Connes' cyclic cohomology and Hochschild cohomology:
\[
HC^n(A) \subseteq H^n_{\mathrm{cyc}}(A, \mathbb{F}) \subseteq HH^n(A, \mathbb{F}).
\]
\end{corollary}

\begin{remark}
For a commutative associative algebra $A$, all three cohomologies coincide up to a shift:
\[
HC^n(A) \cong H^n_{\mathrm{cyc}}(A, \mathbb{F}) \cong HH^n(A, \mathbb{F}) \cong \Omega^n(A) / d\Omega^{n-1}(A),
\]
where $\Omega^n(A)$ are Kähler differentials. In the non-commutative cyclic associative case, the inclusions are proper.
\end{remark}

\section{Appendix}
\subsection*{Detailed Proof of Proposition \ref{prop:cohomology_preserved}}\label{appendix:proof}




Let $f \in C^n_{\mathrm{cyc}}(A, M)$ satisfy conditions (\ref{6})-(\ref{10}). Recall the Hochschild coboundary:
\begin{align*}
(\delta f)(a_0, \dots, a_{n+1}) = a_0 f(a_1, \dots, a_{n+1}) + \sum_{k=0}^{n} (-1)^{k+1} f(a_0, \dots, a_k a_{k+1}, \dots, a_{n+1})\\ + (-1)^{n+1} f(a_0, \dots, a_n) a_{n+1}.
\end{align*}

We will prove that $\delta f$ satisfies each condition.

\medskip
\paragraph{Preservation of Condition (\ref{6}).}
We must show for all $a, a_1, \dots, a_{n+1} \in A$ and any $i$:
\[
a \cdot (\delta f)(a_1, \dots, a_i, \dots, a_{n+1}) = a_i \cdot (\delta f)(a_{i+1}, \dots, a_{n+1}, a, a_1, \dots, a_{i-1}). \qquad (\ref{6})
\]

\subparagraph{Step 1} Write the left-hand side $L = a \cdot (\delta f)(a_1, \dots, a_{n+1})$ as:
\[
L = a \cdot \Bigg[ a_1 f(a_2, \dots, a_{n+1}) + \sum_{k=1}^{n} (-1)^{k+1} f(a_1, \dots, a_k a_{k+1}, \dots, a_{n+1}) + (-1)^{n+1} f(a_1, \dots, a_n) a_{n+1} \Bigg].
\]

\subparagraph{Step 2} Distribute $a \cdot$:
\[
L = \underbrace{a \cdot (a_1 f(a_2, \dots, a_{n+1}))}_{L_1} + \sum_{k=1}^{n} (-1)^{k+1} \underbrace{a \cdot f(a_1, \dots, a_k a_{k+1}, \dots, a_{n+1})}_{L_2(k)} + (-1)^{n+1} \underbrace{a \cdot (f(a_1, \dots, a_n) a_{n+1})}_{L_3}.
\]

\subparagraph{Step 3} Simplify $L_1$ using the module action and (\ref{6}):
\[
L_1 = (a a_1) \cdot f(a_2, \dots, a_{n+1}).
\]
Apply (\ref{6}) to $f$ with external multiplier $a a_1$ and arguments starting at $a_2$:
\[
(a a_1) \cdot f(a_2, \dots, a_i, \dots, a_{n+1}) = a_i \cdot f(a_{i+1}, \dots, a_{n+1}, a a_1, a_2, \dots, a_{i-1}).
\]
Thus
\[
L_1 = a_i \cdot f(a_{i+1}, \dots, a_{n+1}, a a_1, a_2, \dots, a_{i-1}).
\]

\subparagraph{Step 4} Simplify $L_2(k)$. We consider three cases.

\noindent\textit{Case 1: $i < k$.} The $i$-th argument $a_i$ is before the product $a_k a_{k+1}$. Apply (6) to $f$:
\[
a \cdot f(a_1, \dots, a_i, \dots, a_k a_{k+1}, \dots) = a_i \cdot f(a_{i+1}, \dots, a_k a_{k+1}, \dots, a_{n+1}, a, a_1, \dots, a_{i-1}).
\]

\noindent\textit{Case 2: $i > k+1$.} The $i$-th argument $a_i$ is after the product. Apply (\ref{6}) similarly:
\[
a \cdot f(a_1, \dots, a_k a_{k+1}, \dots, a_i, \dots) = a_i \cdot f(a_{i+1}, \dots, a_{n+1}, a, a_1, \dots, a_k a_{k+1}, \dots, a_{i-1}).
\]

\noindent\textit{Case 3: $i = k$ or $i = k+1$.} The distinguished argument is inside the product. Use (\ref{10}):
\[
a \cdot f(\dots, a_k a_{k+1}, \dots) = (a_k a_{k+1}) \cdot f(\dots, a, \dots).
\]
Then apply (\ref{6}) to $(a_k a_{k+1}) \cdot f(\dots)$. The result matches the corresponding term on the right-hand side after a cyclic permutation.

\subparagraph{Step 5} Simplify $L_3$:
\[
L_3 = a \cdot (f(a_1, \dots, a_n) \cdot a_{n+1}) = (a \cdot f(a_1, \dots, a_n)) \cdot a_{n+1}.
\]
Apply (\ref{6}) to $f$ with external multiplier $a$:
\[
a \cdot f(a_1, \dots, a_n) = a_1 \cdot f(a_2, \dots, a_n, a).
\]
Thus
\[
L_3 = (a_1 \cdot f(a_2, \dots, a_n, a)) \cdot a_{n+1} = a_1 \cdot (f(a_2, \dots, a_n, a) \cdot a_{n+1}).
\]

\subparagraph{Step 6} The right-hand side $R = a_i \cdot (\delta f)(a_{i+1}, \dots, a_{n+1}, a, a_1, \dots, a_{i-1})$ expands to:
\begin{align*}
R = a_i \cdot \Bigg[ a_{i+1} f(a_{i+2}, \dots, a_{n+1}, a, a_1, \dots, a_{i-1}) + \sum_{k} (\pm) f(\dots) \\+ (-1)^{n+1} f(a_{i+1}, \dots, a_{n+1}, a, a_1, \dots, a_{i-2}, a_{i-1}) \cdot a_{i-1} \Bigg].
\end{align*}

Comparing term by term, each component of $L$ matches a component of $R$ after applying the cyclic associative identities of $A$. Hence $L = R$.

\medskip
\paragraph{Preservation of Condition (\ref{7}).}
We must show for $1 \le j \le n+1$:
\[
(\delta f)(a_0, \dots, a_{j-1}, b a_j, a_{j+1}, \dots, a_{n+1}) = (\delta f)(a_0, \dots, b, a_j a_{j-1}, a_{j+1}, \dots, a_{n+1}). \quad (\ref{7}) 
\]

\subparagraph{Step 1} Write the left-hand side $L = (\delta f)(a_0, \dots, a_{j-1}, b a_j, a_{j+1}, \dots, a_{n+1})$.

\subparagraph{Step 2} Consider the first term $T_1 = a_0 f(a_1, \dots, a_{j-1}, b a_j, a_{j+1}, \dots, a_{n+1})$. Apply (\ref{7}) to $f$:
\[
f(a_1, \dots, a_{j-1}, b a_j, \dots) = f(a_1, \dots, b, a_j a_{j-1}, \dots).
\]
Thus $T_1$ becomes $a_0 f(a_1, \dots, b, a_j a_{j-1}, \dots)$, which is exactly the first term of the right-hand side.

\subparagraph{Step 3} For the terms $T_2(k)$ with $k \neq j-1$, apply (\ref{7}) directly to $f$ inside. The transformation is straightforward.

\subparagraph{Step 4} For $k = j-1$, we have $f(a_0, \dots, a_{j-2}, a_{j-1}(b a_j), a_{j+1}, \dots)$. By the cyclic associative identity:
\[
a_{j-1}(b a_j) = (a_{j-1} b) a_j = b (a_j a_{j-1}).
\]
Hence
\[
f(a_0, \dots, a_{j-2}, a_{j-1}(b a_j), \dots) = f(a_0, \dots, a_{j-2}, b (a_j a_{j-1}), \dots).
\]
Now apply (7) again to move $b$ to the left:
\[
f(a_0, \dots, a_{j-2}, b (a_j a_{j-1}), \dots) = f(a_0, \dots, b, (a_j a_{j-1}) a_{j-2}, \dots).
\]
This matches the corresponding term on the right-hand side.

\subparagraph{Step 5} The last term $T_3$ transforms similarly using (\ref{7}). Thus $L = R$.

\medskip
\paragraph{Preservation of Condition (\ref{8}).}
We must show:
\[
a \cdot (\delta f)(b a_0, a_1, \dots, a_{n+1}) = b \cdot (\delta f)(a_0 a, a_1, \dots, a_{n+1}).\quad (\ref{8})
\]

\subparagraph{Step 1} Compute the left-hand side $L = a \cdot (\delta f)(b a_0, a_1, \dots, a_{n+1})$.

\subparagraph{Step 2} For $T_1 = a \cdot ((b a_0) f(a_1, \dots, a_{n+1})) = (a(b a_0)) f(a_1, \dots, a_{n+1})$.
Using the cyclic identity $a(b a_0) = b(a_0 a)$, we get $b \cdot (a_0 a f(a_1, \dots, a_{n+1}))$, which is $b$ times the first term of the right-hand side.

\subparagraph{Step 3} For each $T_2(k) = a \cdot f(b a_0, a_1, \dots, a_k a_{k+1}, \dots, a_{n+1})$, apply (\ref{8}) to $f$:
\[
a \cdot f(b a_0, a_1, \dots) = b \cdot f(a_0 a, a_1, \dots).
\]
Thus $T_2(k) = b \cdot f(a_0 a, a_1, \dots, a_k a_{k+1}, \dots)$, which is $b$ times the corresponding $T_2(k)$ term on the right-hand side.

\subparagraph{Step 4} For $T_3 = a \cdot (f(b a_0, a_1, \dots, a_n) a_{n+1}) = (a \cdot f(b a_0, a_1, \dots, a_n)) a_{n+1}$.
Apply (\ref{8}) to $f$: $a \cdot f(b a_0, a_1, \dots, a_n) = b \cdot f(a_0 a, a_1, \dots, a_n)$.
Hence $T_3 = (b \cdot f(a_0 a, a_1, \dots, a_n)) a_{n+1} = b \cdot (f(a_0 a, a_1, \dots, a_n) a_{n+1})$, which is $b$ times the last term of the right-hand side.

\subparagraph{Step 5} Summing over all terms gives $L = b \cdot (\delta f)(a_0 a, a_1, \dots, a_{n+1})$.

\medskip
\paragraph{Preservation of Condition (\ref{9}).}
We must show for $0 \le j \le n+1$:
\[
a \cdot a'_j \cdot (\delta f)(a_0, \dots, a_j, \dots, a_{n+1}) = a'_j \cdot a_j \cdot (\delta f)(a_0, \dots, a, \dots, a_{n+1}). \qquad (\ref{9})
\]

\subparagraph{Step 1} Apply $\delta f$ to the left side. Each term is of the form $a \cdot a'_j \cdot (a_0 \cdot f(\dots))$ or $a \cdot a'_j \cdot (f(\dots) \cdot a_{n+1})$.

\subparagraph{Step 2} Using the cyclic bimodule axioms:
\[
a \cdot a'_j \cdot (a_0 \cdot m) = (a a_0) \cdot (a'_j \cdot m).
\]
Now apply (\ref{9}) to $f$:
\[
(a a_0) \cdot (a'_j \cdot f(\dots, a_j, \dots)) = (a a_0) \cdot (a'_j a_j \cdot f(\dots, a, \dots)) = a'_j a_j \cdot ((a a_0) \cdot f(\dots, a, \dots)).
\]

\subparagraph{Step 3} The other terms $f(\dots) \cdot a_{n+1}$ are handled similarly using the module axioms and (\ref{9}). Summing over all terms yields the right-hand side.

\medskip
\paragraph{Preservation of Condition (\ref{10}).}
Condition (\ref{10}) is a special case of (\ref{6}) where the $i$-th argument $a_i$ is replaced by $a_i b$. Since we have already proved that (\ref{6}) is preserved by $\delta f$, and the substitution $a_i \mapsto a_i b$ does not affect the linearity of $\delta$, condition (\ref{10}) holds for $\delta f$.

\medskip
Thus $\delta(f)$ satisfies all conditions (\ref{6})-(\ref{10}). 
\section*{Declarations}


\begin{itemize}
\item Funding: Not applicable.
\item Conflict of interest/Competing interests: Not applicable.
\item Ethics approval and consent to participate: Not applicable.
\item Consent for publication: Not applicable.
\item Data availability: Not applicable. 
\item Materials availability: Not applicable.
\item Code availability: Not applicable 
\item Author contribution: Not applicable.
\end{itemize}

\begin{thebibliography}{99}

\bibitem{Abdelwahab2025}
H. Abdelwahab, I. Kaygorodov, B. Sartayev, \textit{Shift associative algebras}, 
Rend. Circ. Mat. Palermo, II. Ser. \textbf{74}, 145 (2025).


\bibitem{Dotsenko2023}
V. Dotsenko, N. Ismailov, U. Umirbaev, \textit{Polynomial identities in Novikov algebras}, Mathematische Zeitschrift 303(60), 3 (2023).

\bibitem{Sartayev2023}
B. Sartayev, P. Kolesnikov, \textit{Noncommutative Novikov algebras}, Eur. J. Math. 9(2), 35 (2023).

\bibitem{Drensky2024}
V. Drensky, N. Ismailov, M. Mustafa, B. Zhakhayev, \textit{Free bicommutative superalgebras}, J. Algebra 652, 158-187 (2024).

\bibitem{Ismailov2024}
N.A. Ismailov, F.A. Mashurov, B.K. Sartayev, \textit{On algebras embeddable into bicommutative algebras}, Comm. Algebra 52(11), 4778-4785 (2024).

\bibitem{Kaygorodov2024}
I. Kaygorodov, F. Mashurov, \textit{Mutations of perm algebras}, Revista de la Real Academia de Ciencias Exactas, Físicas y Naturales. Serie A. Matemáticas 118, 166 (2024).

\bibitem{Mashurov2024}
F. Mashurov, B. Sartayev, \textit{Metabelian Lie and perm algebras}, J. Algebra Appl. 23(4), 2450065 (2024).

\bibitem{Hochschild1945}
G. Hochschild, \textit{On the cohomology groups of an associative algebra}, Ann. Math. 46(1), 58-67 (1945).

\bibitem{Connes1985}
A. Connes, \textit{Noncommutative differential geometry}, Publ. Math. Inst. Hautes Études Sci. 62, 41-144 (1985).
\bibitem{CuntzQuillen}
J. Cuntz and D. Quillen, \textit{Algebra extensions and nonsingularity}, 
J. Amer. Math. Soc. \textbf{8} (1995), no. 2, 251--289.
\end{thebibliography}

%% if required, the content of .bbl file can be included here once bbl is generated
%%\input sn-article.bbl

\end{document}
